\section{Method}
\label{sec:method}

\subsection{Views of an item}

An evaluation item is a context whose span is expressed four ways. Let an item
consist of a prefix $p$, a span $s$, and a suffix $q$, where $p$ and $q$ are
shared verbatim across all four views and only the span differs:
\begin{align}
    \VT   &= p \,\Vert\, s \,\Vert\, q,          &\quad
    \VT'  &= p \,\Vert\, \tilde{s} \,\Vert\, q, \\
    \VI   &= p \,\Vert\, \mathcal{R}(s) \,\Vert\, q, &\quad
    \VI'  &= p \,\Vert\, \mathcal{R}'(s) \,\Vert\, q,
\end{align}
where $\mathcal{R}$ renders the span into the visual channel and $\tilde{s}$ is
a same-modality re-expression of the same content. The primed views are the
\emph{within-modality controls}: they hold content fixed and vary only surface
form, and they exist to set the scale against which the cross-modal distance is
read.

The suffix must be non-empty. It is where the comparison is made, and an empty
suffix would leave no position at which the two views can be required to agree.

\subsection{The merge position}
\label{sec:merge}

$\VT$ and $\VI$ have different lengths: in our configuration a one-word span
costs one to two text tokens and six visual tokens. There is therefore no
token-level correspondence along the span, and we assert nothing about it.
Instead we compare the hidden state at the first token of the shared suffix,
\begin{equation}
    m = \text{len}(\text{sequence}) - \text{len}(\text{suffix tokens}),
    \label{eq:merge}
\end{equation}
computed per item from the attention mask so that it is correct under padding.
Two properties make this the right choice. It is model-agnostic --- it never
requires knowing how the processor expanded the image --- and it is the earliest
position at which the two views are required to be comparable, since from there
on they are conditioned on identical text.

Equation~\ref{eq:merge} assumes the suffix tokenises identically in both views.
Tokenisation is not compositional, so a merge across the span/suffix boundary
would silently shift every index by one and corrupt every downstream number
without raising. We therefore verify at run time that the suffix is a token-level
suffix of each sequence, and abort otherwise. The same non-compositionality
governs how we count option lengths in \S\ref{sec:validity}: token counts are
always obtained by differencing sequences encoded through the processor, never
by encoding a fragment in isolation.

Forward passes use right padding so that absolute indices computed on the
unpadded sequence remain valid; batched generation, used only for the
transcription check, requires left padding instead.

\subsection{The Modality Substitution Gap}
\label{sec:msg}

Let $h(\cdot)$ denote the hidden state at the merge position at a chosen layer,
and $d(\cdot,\cdot)$ the cosine distance. The normalised gap for an item is
\begin{equation}
    \msg = \frac{d\big(h(\VT),\, h(\VI)\big)}
                {\tfrac{1}{2}\Big[d\big(h(\VT),\, h(\VT')\big)
                 + d\big(h(\VI),\, h(\VI')\big)\Big]}.
    \label{eq:msg}
\end{equation}
The denominator is what makes this interpretable. It asks how far apart two
expressions of the \emph{same} content sit within a single modality, on the same
items and at the same position, and the ratio then has a meaningful reference
value: $\msg \approx 1$ means crossing modalities costs no more than
re-expressing content within one, $\msg \gg 1$ means the channel dominates the
content, and $\msg < 1$ would mean surface variation within a modality is the
larger effect.

We report the ratio of means, $\overline{d_{\text{cross}}}/\overline{d_{\text{within}}}$,
as the headline, with the mean of per-item ratios alongside. The two answer
different questions and can diverge sharply when the denominator is
small for some items; per-item ratios are heavy-tailed for exactly that reason,
which is why the ratio of means is primary. Confidence intervals are obtained by
bootstrap over items. We additionally report the two halves of the denominator
separately, since a denominator dominated by one modality's control is really
being set by that control alone.

\paragraph{The denominator is a design decision, not a detail.}
A within-modality control that is too weak inflates $\msg$. A capitalisation
variant preserves content but perturbs the representation very little, so a gap
measured against it is a statement about a small reference quantity; a synonym is
a genuinely different expression of the same content and is the stronger
comparison. We report which control was used with every number, because
``$\msg = 3$'' means different things against the two.

\subsection{Three substitution tiers}
\label{sec:tiers}

\paragraph{Tier A: referential.} A noun phrase replaced by a crop of the region
it refers to, using existing phrase-to-region annotations. The bridge is
grounding. This tier is a \emph{control}, not the contribution: the crop and the
phrase genuinely differ in content, so any gap here is an upper bound mixing
representational divergence with information asymmetry. The image control is a
different region annotated with the same phrase.

\paragraph{Tier B: orthographic.} A word replaced by a rendering of that word as
glyphs. This is the centrepiece, for two reasons. The substitution is exactly
content-preserving --- the same lexical item through a different encoder --- so
information asymmetry cannot explain a gap. And it lifts the restriction to
depictable content: abstract words, function words, and rare words render as
readily as concrete nouns, so the tier supports a word-class comparison that
grounding-based probes cannot. The image control is the same word in a different
typeface, which is the exact analogue of a synonym in the text view: same
content, different surface.

Tier B depends on the model being able to read the rendering, so this is
established before anything else is measured, by a transcription check on
held-out typefaces; the rendering configuration is frozen at the cheapest
setting that passes, and every subsequent number is reported both
unconditionally and conditioned on spans the model transcribes correctly. This
conditioning is what keeps a residual transcription error rate from acting as a
confound, and it matters for the word-class comparison specifically: if
transcription accuracy varies with word length or class, an unconditioned
comparison would report legibility as though it were representational structure.

\paragraph{Tier C: relational.} A described relation replaced by a diagram of
it, probing whether relational structure survives substitution. This tier is the
most demanding on stimulus design, and \S\ref{sec:validity} explains why it
cannot be trusted without an explicit comprehension check.

\subsection{Validity: what the gap is not}
\label{sec:validity}

$\msg$ is defined whether or not the substituted span means anything to the
model, and will report a large, tight-CI gap for a stimulus the model cannot
read. Two families of control are therefore not optional.

\paragraph{Does the image view carry the content?}
We score a two-alternative forced choice --- the true span against a distractor
drawn from the same class --- under each view, by likelihood rather than by
parsing generated text, which avoids refusals, formatting drift, and any
prompt-following confound between modalities. Options are scored as pointwise
mutual information against a null context,
\begin{equation}
    \text{score}(o) = \log P(o \mid \text{context}) - \log P(o \mid \text{question alone}),
\end{equation}
because raw likelihood lets an option's unconditional frequency compete with the
in-context evidence: a rare true span can lose to a common distractor even with
the answer written out. The null term is computed once per option string and
shared across views, so the calibration cannot differ by modality.

Accuracy on the image view is the validity check, but it must be read against
the right reference. Chance is not it: distractors from the same class in a
constraining context are partly predictable without the span. We therefore score
a third, \emph{span-free} view with the span blanked in both modalities, and
require the image view to clear that floor. A choice decidable from the context
alone is reported as such, and the tier's $\msg$ is marked uninterpretable.

A saturated view is likewise uninformative in the other direction: an image view
at ceiling has no headroom, so a text-image accuracy difference measured against
it bounds the functional gap near zero rather than measuring it. We report that
condition explicitly rather than quoting the difference.

\paragraph{Is the gap a difference in information, or in coordinates?}
A representational difference that any linear readout could undo costs the model
nothing. This matters concretely, because the projector bridging a VLM's vision
encoder to its language model is itself a linear map: a gap that a fitted linear
map removes is one the standard architecture already closes. We therefore
recompute $\msg$ after removing, in turn,
\begin{enumerate}
    \item a constant per-modality translation (the offset-free gap),
    \item an orthogonal map, fitted by Procrustes,
    \item an arbitrary linear map, fitted by ridge regression,
\end{enumerate}
each strictly more general than the last. The level at which the gap disappears
is the finding.

Three details decide whether this analysis means anything. \emph{The maps are
fitted out-of-fold}: fitted and scored on the same items, a $[D,D]$ map has
enough freedom to drive any distance to zero. \emph{The within-modality control
is carried through the same fold's map} as the cross-modal pair, since mapping
the numerator while leaving the denominator alone would manufacture a collapse
in the ratio. And \emph{the fit must not be under-determined}: a map with $D^2$
parameters fitted from fewer than $D$ rows generalises poorly whatever the truth
is, and that failure is indistinguishable from an irreducible gap. We report
rows-per-dimension with every fitted result and decline to issue a verdict below
a fixed threshold. The ridge penalty is selected from a grid by out-of-fold
distance --- mildly optimistic for the map, deliberately, since the claim under
test is that the gap \emph{survives} it.

\paragraph{Informativeness.}
Invariance is trivially maximised by discarding information, so a low $\msg$ is
only meaningful alongside evidence that the representation still distinguishes
spans. We train a linear probe to recover span identity from the merge-position
state and report probe accuracy alongside $\msg$ throughout, treating the pair as
a frontier rather than either as a scalar objective.

\subsection{Complementary measurements}

Beyond the merge-position geometry we report: the Jensen--Shannon divergence
between the two views' next-token distributions over the shared suffix, which
asks whether the residual geometric difference is behaviourally visible; linear
CKA between the two views' representations, which asks whether the two clouds
have the same \emph{shape} independent of orientation and is therefore the
natural companion to the null hierarchy; per-layer distances, which locate where
in the stack the gap lives; and a text$\rightarrow$image$\rightarrow$text
transcription cycle, which bounds how much of the span survives a round trip
through the visual channel.

\subsection{Experimental protocol}

All gated measurements are taken on a single GPU architecture with a fixed
software stack, and every recorded run carries its architecture, framework
version, and attention backend, since a comparison split across architectures is
not a comparison. Decision thresholds and their consequences are fixed in a
version-controlled document before the corresponding numbers exist, and each
decision is recorded on its date with the numbers available on that date. We
state this as method rather than as good practice because two of our own
intermediate conclusions were reversed by controls added afterwards, and the
record of what was decided when is the only thing that distinguishes a
pre-registered threshold from a post-hoc one.
